% Desarrollo con asistencia de IA; pendiente de revisión del integrante responsable.
\subsection{Contador lógico por proceso}
El componente \path{src/common/lamport.py} define la clase
\texttt{LamportClock}. Cada nodo debe tener su propia instancia, compartida
entre los hilos de servidor y cliente del mismo proceso. No se comparte un
contador entre los tres nodos: eso eliminaría la independencia de sus relojes.
Dos procesos separados tampoco comparten esta instancia aunque utilicen el
mismo nombre de nodo.

Las operaciones de actualización y consulta usan un bloqueo de exclusión mutua.
El método \texttt{tick()} incrementa y devuelve el contador; \texttt{send()}
es una interfaz equivalente para envíos. La recepción valida primero que la
marca sea un entero no negativo y después aplica la regla de
la ecuación~\eqref{eq:lamport-recepcion}. Se rechazan valores booleanos aunque
Python los considere subtipos de entero. La propiedad \texttt{value} permite
consultar el reloj sin modificarlo.

El siguiente fragmento muestra la sección crítica de recepción:
\begin{lstlisting}[language=Python]
def receive(self, remote: int) -> int:
    _validate_timestamp(remote)
    with self._lock:
        self._value = max(self._value, remote) + 1
        return self._value
\end{lstlisting}

Para conservar compatibilidad de interfaces, el módulo también proporciona
\texttt{enviar()}, \texttt{recibir()} y \texttt{valor}. Los alias no representan
eventos adicionales: se elige una sola operación por evento. La compatibilidad
de nombres no demuestra por sí sola la integración con un transporte.

El cliente y el servidor TCP de esta revisión importan la clase común, en lugar
de mantener implementaciones locales duplicadas. Conservan sus argumentos de
ejecución y el formato JSON existente. El servidor y el cliente gRPC revisados
también importan la clase común y el servidor incrementa antes de enviar el acuse.
Esto unifica el algoritmo, pero los ejecutables actuales todavía crean instancias
separadas para cliente y servidor. No constituyen por sí solos una arquitectura
de tres procesos con servidor y cliente compartiendo reloj dentro de cada nodo.
El registro TCP existente también protege con un bloqueo la creación de cabecera
y la escritura de filas entre hilos del proceso. Este bloqueo evita escrituras
simultáneas, pero no garantiza que el orden de las filas coincida con el orden
de actualización de un reloj protegido por otro bloqueo independiente.

\subsection{Registro e identificación de eventos}
El componente propuesto en \path{src/common/event_log.py} implementa
\texttt{EventLog}. Su bloqueo engloba la actualización del reloj y la escritura
del evento, de modo que los hilos del proceso no registren una marca obtenida
por otra operación concurrente. El método devuelve exactamente la marca
registrada; ese retorno debe colocarse en el mensaje enviado.

\begin{lstlisting}[language=Python]
timestamp = log.record(
    "envio_solicitud", message_id, origin, destination)
# El transporte envia timestamp y message_id.
received = log.record(
    "recepcion_solicitud", message_id, origin,
    destination, remote=request_timestamp)
\end{lstlisting}

Como \texttt{record()} ya actualiza el contador, no se llama además a
\texttt{send()} o \texttt{receive()} para ese mismo evento. Durante una captura
completa, todos los eventos del reloj deben pasar por el registro y la captura
debe comenzar con el contador en cero.

El CSV contiene mecanismo, corrida, nodo, secuencia local, tipo de evento,
reloj local, reloj remoto cuando corresponde, hora UTC, identificador de
mensaje, origen y destino. La solicitud y su acuse comparten identificador,
pero sus tipos de evento permiten distinguir las dos direcciones. El ID debe
ser único por intercambio dentro de cada mecanismo y corrida y debe
transmitirse hasta el receptor; no se reconstruye a partir del orden de llegada.

Cada proceso escribe un archivo nuevo y exclusivo por corrida. La apertura
rechaza archivos existentes para evitar mezclar ejecuciones. La secuencia local
establece el orden del proceso; la hora UTC documenta cuándo se registró el
evento y no se utiliza como prueba de causalidad.

Este esquema corresponde al componente de registro preparado, no a los CSV
que producen actualmente los adaptadores del equipo. Sus registros existentes
no incluyen corrida, secuencia ni reloj remoto; el servidor gRPC reconstruye
los identificadores mediante contadores por emisor y conexión. Para emplear el
verificador completo se necesita integrar el registro y transmitir el ID real
desde el cliente. No es válido completar los CSV antiguos con valores supuestos.

\subsection{Concurrencia, validación y límites}
Las pruebas automatizadas comprueban actualizaciones sin pérdida entre hilos,
filas completas y secuencias ordenadas. El verificador
\path{analysis/verify_lamport.py} comprueba las reglas del contador y empareja
envíos y recepciones de solicitudes y acuses. También rechaza secuencias
incompletas, duplicados y marcas transmitidas inconsistentes.

El bloqueo no sincroniza procesos distintos ni protege automáticamente una
conexión de red compartida. Los adaptadores TCP y gRPC deben gestionar su
transporte y transmitir las marcas e identificadores del registro. Un fallo
de escritura invalida la captura y debe interrumpir esa corrida.

Un timeout no demuestra que el receptor no haya procesado la solicitud. Un
reintento no debe añadirse como si fuera el mismo intercambio completo: requiere
una política explícita de identificación y, si se pretende evitar efectos
duplicados, deduplicación en el receptor. El componente de reloj no implementa
esa política ni detecta fallos de nodos.

\pendiente{Integrar estos componentes en ambos transportes, ejecutar las pruebas
de tres clientes y fallos de nodos, y documentar los resultados reales. Las
pruebas unitarias no acreditan todavía esas ejecuciones.}
